\documentclass[11pt]{article}

\usepackage[margin=1in]{geometry}
\usepackage{amsmath,amssymb}
\usepackage{booktabs}
\usepackage{microtype}
\usepackage[hidelinks]{hyperref}

\newcommand{\Hil}{\mathcal{H}}
\newcommand{\D}{D}

\title{The Hilbert-Transform Surface Channel in CS-DNO\\
Rationale, Evidence Gap, and Deletion Ablation}
\author{}
\date{August 22, 2026}

\begin{document}
\maketitle

\section{Question and present status}

The C27 CS-DNO represents the surface with the fixed feature stack
\begin{equation}
E_0(\eta)
=
\left[
\eta,
\eta^2,
\eta^3,
\partial_x\eta,
\partial_x^2\eta,
|\D|^{1/2}\eta,
\Hil\eta
\right].
\label{eq:feature-stack}
\end{equation}
The final channel was inherited from the early broad CS-DNO feature bank.
It was never isolated in an ablation.  Its inclusion is therefore a plausible
architectural hypothesis, not an experimentally established requirement of
C27.

\section{What transform is used?}

For a real periodic surface, the spatial Hilbert transform is the Fourier
multiplier
\begin{equation}
\widehat{\Hil\eta}(k)
=
-\mathrm{i}\,\operatorname{sgn}(k)\,\widehat{\eta}(k).
\label{eq:hilbert-symbol}
\end{equation}
It preserves the magnitude of every nonzero Fourier mode and shifts its phase
by ninety degrees.  It is an order-zero, nonlocal operator: unlike
$\partial_x\eta$ or $\partial_x^2\eta$, it does not amplify high wavenumbers.
It also satisfies
\begin{equation}
\Hil\partial_x = |\D|,
\label{eq:hilbert-derivative}
\end{equation}
which connects it to the nonlocal multipliers that occur in water-wave
operators.

This channel is unrelated to the strip-Hilbert transform used in constructing
Tanaka initial conditions.  Here it is only a deterministic input feature of
the learned surface trunk.

\section{Why was it plausible?}

After the fixed feature construction, the shared surface trunk is pointwise in
$x$.  Pointwise powers of $\eta$ cannot reconstruct the quadrature component
$\Hil\eta$.  Supplying it explicitly gives the trunk access to nonlocal phase
information without asking a learned local map to synthesize a Fourier
operator.  A concise explanation, if the channel survives ablation, is:

\begin{quote}
We include the periodic Hilbert transform of the surface to expose nonlocal
quadrature information to an otherwise pointwise surface-feature trunk.  The
channel is order zero, so it supplies phase information without the
high-wavenumber amplification of derivative channels.
\end{quote}

This explains what information the channel makes available.  It does not show
that the information improves the learned DNO.

\section{What evidence do we actually have?}

The historical minimal-feature experiment compared the complete feature bank
against a model that simultaneously reduced the polynomial degree and disabled
all spectral surface channels.  The deletion regressed validation loss by a
factor of approximately $2.3$, establishing that the \emph{feature bank as a
whole} matters.  It does not identify the contribution of $\Hil\eta$.

In particular:
\begin{itemize}
  \item no run changed only the Hilbert channel;
  \item neither the analytic $G_0+G_1$ backbone nor the three C27 loss terms
        requires an explicit $\Hil\eta$ input;
  \item the current evidence cannot distinguish a useful phase feature from a
        redundant inherited channel.
\end{itemize}

The defensible pre-ablation statement is therefore that $\Hil\eta$ is an
unisolated, physics-motivated feature.  It should not be described as
load-bearing on the basis of the existing grouped ablation.

\section{Hilbert-only deletion protocol}

Run a one-variable comparison after fixing the full-row C27 optimizer-step
budget.  Both arms must use the same fresh seed, authenticated corpus, row
order, normalization, batch size, optimizer schedule, and three loss terms.

\begin{center}
\begin{tabular}{lll}
\toprule
Arm & Hilbert channel & All other settings \\
\midrule
Control & enabled & locked C27 recipe \\
Deletion & disabled & identical to control \\
\bottomrule
\end{tabular}
\end{center}

The trainer already exposes the required switch:
\begin{center}
\texttt{--cs\_use\_hilbert}
\qquad versus \qquad
\texttt{--no-cs\_use\_hilbert}.
\end{center}

At the locked C27 dimensions, the control has $1{,}342{,}400$ parameters and
the deletion has $1{,}342{,}080$.  The difference is only $320$ parameters,
approximately $0.024\%$, so the comparison is effectively capacity matched.

Evaluate both held-out one-step error and the frozen case-by-case rollout
panel.  The rollout comparison must include nonfinite counts, raw and
translation-aligned tail errors, and per-family worst cases.  A validation-only
win is insufficient if rollout stability regresses.

\section{Decision rule}

Retain $\Hil\eta$ only if the control shows a reproducible accuracy or stability
benefit larger than run-to-run variation.  If the deletion is equivalent or
better, remove the channel and prefer the smaller, easier-to-defend feature
stack.  Until that comparison is complete, preserve C27 when reproducing the
locked model, but mark the channel as unresolved rather than justified.

\end{document}
